% Work-MAT | 2. rocnik | hodina 18
\documentclass[11pt,a4paper]{article}
\input{../../sablony/pripravy/priprava.sty}
\usepackage{tikz}
\usetikzlibrary{arrows.meta}
\nastavPripravu{2. ročník}{Exponenciální rovnice}{18}{Exponenciální rovnice VI: složitější substituce}
\begin{document}
\poznamka{Postup od známého typu $t^2+bt+c=0$ k podílům a záporným exponentům. Bez logaritmů. Orientačně 6 min návrat, 15 min nový trik, 13 min procvičení, 11 min diskuse o volbě metody.}
\blok{Když v rovnici vystupuje $a^x$ i $a^{-x}$}
Díky pravidlům pro mocniny můžeme vyjádřit $a^{-x}=1/a^x$. Proto při substituci $t=a^x>0$ dostaneme $a^{-x}=1/t$. Po vynásobení nenulovým $t$ se někdy opět objeví kvadratická rovnice.
\begin{prikladReseniBox}[title={Příklad 1}]
Vyřeš $2^x+2^{-x}=\frac52$.
\tcblower
\[\begin{aligned}t&=2^x>0,\\t+\frac1t&=\frac52,\\2t^2+2&=5t,\\2t^2-5t+2&=0,\\(2t-1)(t-2)&=0.\end{aligned}\]
Oba pomocné kořeny jsou kladné: $t_1=\frac12$, $t_2=2$. Návrat:
\[2^x=2^{-1}\Rightarrow x=-1,\qquad2^x=2^1\Rightarrow x=1.\]
$K=\{-1,1\}$; zkouška v obou případech dává $2+\frac12=\frac52$.
\end{prikladReseniBox}
\begin{prikladReseniBox}[title={Příklad 2}]
Vyřeš $2^x+2^{3-x}=6$.
\tcblower
Položme $t=2^x>0$. Pak $2^{3-x}=8/2^x=8/t$. Dostáváme
\[\begin{aligned}t+\frac8t&=6,\\t^2-6t+8&=0,\\(t-2)(t-4)&=0.\end{aligned}\]
Z $t=2$ plyne $x=1$, z $t=4$ plyne $x=2$. Zkouška: $2+4=6$; $4+2=6$. $K=\{1,2\}$.
\end{prikladReseniBox}
\blok{Po sjednocení základu může být substituce jednoduchá}
\begin{prikladReseniBox}[title={Příklad 3}]
Vyřeš $2\cdot4^x+4\cdot2^x=48$.
\tcblower
\[\begin{aligned}t&=2^x>0,\\2t^2+4t&=48,\\t^2+2t-24&=0,\\(t+6)(t-4)&=0.\end{aligned}\]
Přijmeme pouze $t=4$, tedy $2^x=4$ a $x=2$. Zkouška: $2\cdot16+4\cdot4=48$.
\end{prikladReseniBox}
\poznamka{První dvě úlohy vyžadují násobení $t$. Zdůrazni, že to smíme provést, protože $t>0$, takže nedělíme ani nenásobíme nulou. Záporný kořen pomocné rovnice se vždy vylučuje.}
\begin{zadaniReseniBox}[title={Samostatná práce 1}]
\begin{enumerate}[label=\alph*)]
\item $3^x+3^{2-x}=6$;
\item $2^x+2^{2-x}=5$;
\item $2\cdot4^x-6\cdot2^x+4=0$.
\end{enumerate}
\tcblower
\begin{enumerate}[label=\alph*)]
\item $t+9/t=6$, tj. $(t-3)^2=0$, $t=3$, tedy $x=1$.
\item $t+4/t=5$, tj. $(t-1)(t-4)=0$, tedy $x=0$ nebo $x=2$.
\item $2t^2-6t+4=0$, tj. $(t-1)(t-2)=0$, tedy $x=0$ nebo $x=1$.
\end{enumerate}
\end{zadaniReseniBox}
\begin{zadaniReseniBox}[title={Rozšíření 1}]
Vyřeš $3^{x}+3^{-x}=\frac{10}{3}$.
\tcblower
$t=3^x>0$. Z $t+1/t=10/3$ dostaneme $3t^2-10t+3=(3t-1)(t-3)=0$. Pro $t=1/3,3$ plyne $K=\{-1,1\}$.
\end{zadaniReseniBox}
\begin{zadaniReseniBox}[title={Domácí úkol}]
\begin{enumerate}[label=\alph*)]
\item Vyřeš $2^x+2^{1-x}=3$.
\item Vyřeš $4^x-5\cdot2^x+4=0$ a popiš volbu substituce.
\item Vysvětli, proč je vždy povoleno násobit $t$ v rovnici s $t=5^x$.
\end{enumerate}
\tcblower
\begin{enumerate}[label=\alph*)]
\item $t+2/t=3$, tj. $(t-1)(t-2)=0$, tedy $x=0,1$.
\item $t=2^x>0$, $(t-1)(t-4)=0$, tedy $x=0,2$.
\item Protože $t=5^x>0$, takže $t\ne0$.
\end{enumerate}
\end{zadaniReseniBox}
\end{document}
